% from ega3-6-fr.tex
\oldpage[III]{155}
On voit donc que ${}^{(b)}\mathscr E_{pq}^2$ est le $(-p)$-ème module de
cohomologie du complexe
\[
  L^\bullet(\Phi^{(1)},\Phi^{(2)};\mathscr S)
\]
des cochaînes bi-alternées sur $\Phi^{(1)}$ et $\Phi^{(2)}$ à valeurs dans le
système de coefficients
\[
  \mathscr S:(\sigma,\tau)\leadsto
  \Gamma\bigl(U_\sigma^{(1)}\times_S U_\tau^{(2)},
    \mathscr{T}\!or_q^S
    (\mathscr P_\bullet^{(1)},\mathscr P_\bullet^{(2)})\bigr)
\]
($0_{\mathrm I}$, 11.8.4). On sait alors ($0_{\mathrm I}$, 11.8.5 et 11.8.6)
que la cohomologie de ce complexe est la même que celle du complexe
$C^\bullet(\Phi^{(1)},\Phi^{(2)};\mathscr S)$ de toutes les cochaînes sur
$\Phi^{(1)}$ et $\Phi^{(2)}$ à valeurs dans $\mathscr S$, et aussi que celle du
complexe $P^\bullet(\Phi^{(1)},\Phi^{(2)};\mathscr S)$, dont les éléments sont
les combinaisons linéaires des
\[
  \lambda(\sigma,\tau)\in
  \Gamma\bigl(U_\sigma^{(1)}\times_S U_\tau^{(2)},
    \mathscr{T}\!or_q^S
    (\mathscr P_\bullet^{(1)},\mathscr P_\bullet^{(2)})\bigr),
\]
où $\sigma=(\alpha_0,\ldots,\alpha_h)$ et
$\tau=(\beta_0,\ldots,\beta_h)$ sont des suites ayant même nombre d'éléments.
Mais on a alors
\[
  U_\sigma^{(1)}\times_S U_\tau^{(2)}
  =(U_{\alpha_0}^{(1)}\times_S U_{\beta_0}^{(2)})\cap\cdots\cap
    (U_{\alpha_h}^{(1)}\times_S U_{\beta_h}^{(2)})
  \qquad\text{(I, 3.2.7)}.
\]
Si l'on désigne par $\mathfrak U$ le recouvrement de
$Z=X^{(1)}\times_S X^{(2)}$ par les ouverts affines
$U_\alpha^{(1)}\times_S U_\beta^{(2)}$, on voit finalement, compte tenu de ce
que $X^{(1)}\times_S X^{(2)}$ est un schéma, que l'on a, en vertu de (1.3.1),
\[
  {}^{(b)}\mathscr E_{pq}^2
  =\mathbf H^{-p}\bigl(X^{(1)}\times_S X^{(2)},
    \mathscr{T}\!or_q^S
    (\mathscr P_\bullet^{(1)},\mathscr P_\bullet^{(2)})\bigr).
\]

En second lieu,
$\mathbf{Tor}_q^A(L_{\bullet\bullet}^{(1)},L_{\bullet\bullet}^{(2)})$ est un
bicomplexe dont le terme de degré $(h,k)$ est la somme directe des $A$-modules
\[
  \mathbf{Tor}_q^A
  \bigl(C^{-h_1}(\mathfrak U^{(1)},\mathscr P_{k_1}^{(1)}),
        C^{-h_2}(\mathfrak U^{(2)},\mathscr P_{k_2}^{(2)})\bigr)
\]
tels que $h_1+h_2=h$ et $k_1+k_2=k$; explicitant les modules
$C^r(\mathfrak U^{(i)},\mathscr P_s^{(i)})$ comme ci-dessus, on voit encore
que ce terme est somme directe des $A$-modules
\[
  \Gamma\bigl(U_\sigma^{(1)}\times_S U_\tau^{(2)},
    \mathscr{T}\!or_q^S
    (\mathscr P_{k_1}^{(1)},\mathscr P_{k_2}^{(2)})\bigr),
\]
où $k_1+k_2=k$, et $\sigma$ (resp. $\tau$) parcourt les éléments de
$\mathfrak P(\Phi^{(1)})$ (resp. $\mathfrak P(\Phi^{(2)})$) tels que
$\operatorname{Card}(\sigma)+\operatorname{Card}(\tau)=-h-2$. Le terme
${}^{(b')}\mathscr E_{pq}^2$ que nous calculons est le $(-p)$-ème module de
cohomologie d'un bicomplexe $N^{\bullet\bullet}=(N^{hk})$, où le complexe
simple $N^{\bullet k}$ est le complexe de cochaînes bi-alternées sur
$\Phi^{(1)}$ et $\Phi^{(2)}$, à valeurs dans le système de coefficients
\[
  \mathscr S_k:(\sigma,\tau)\leadsto
  \Gamma\left(U_\sigma^{(1)}\times_S U_\tau^{(2)},
    \bigoplus_{k_1+k_2=k}
    \mathscr{T}\!or_q^S
    (\mathscr P_{-k_1}^{(1)},\mathscr P_{-k_2}^{(2)})\right),
\]
ces systèmes de coefficients formant un complexe $\mathscr S^\bullet$ où la
différentielle provient de celle du complexe simple associé au bicomplexe
$\mathscr{T}\!or_q^S
(\mathscr P_\bullet^{(1)},\mathscr P_\bullet^{(2)})$. On sait que la
cohomologie de $N^{\bullet\bullet}$ est la même que celle du bicomplexe
$C^\bullet(\Phi^{(1)},\Phi^{(2)};\mathscr S^\bullet)$
($0_{\mathrm I}$, 11.8.9), et aussi la même que celle du bicomplexe
$P^\bullet(\Phi^{(1)},\Phi^{(2)};\mathscr S^\bullet)$, où les éléments de
degrés $(h,k)$ sont les combinaisons linéaires des
\[
  \lambda(\sigma,\tau)\in
  \Gamma\left(U_\sigma^{(1)}\times_S U_\tau^{(2)},
    \bigoplus_{k_1+k_2=k}
    \mathscr{T}\!or_q^S
    (\mathscr P_{-k_1}^{(1)},\mathscr P_{-k_2}^{(2)})\right),
\]
$\sigma=(\alpha_0,\ldots,\alpha_h)$,
$\tau=(\beta_0,\ldots,\beta_h)$ étant des suites ayant même nombre
d'éléments ($0_{\mathrm I}$, 11.8.10). On voit alors comme ci-dessus que
${}^{(b')}\mathscr E_{pq}^2$ est le $(-p)$-ème module de cohomologie du
bicomplexe $C^\bullet(\mathfrak U,\mathscr Q^\bullet)$, où
$\mathscr Q^\bullet$ est le complexe simple associé au bicomplexe
$\mathscr{T}\!or_q^S
(\mathscr P_\bullet^{(1)},\mathscr P_\bullet^{(2)})$ de
$\mathscr O_Z$-Modules. Avec les conventions faites dans (6.6.1), on a donc
\[
  {}^{(b')}\mathscr E_{pq}^2
  =\mathbf H^{-p}\bigl(X^{(1)}\times_S X^{(2)},
    \mathscr{T}\!or_q^S
    (\mathscr P_\bullet^{(1)},\mathscr P_\bullet^{(2)})\bigr).
\]
\end{env}

